% ===========================================================================
%  Theorem system - boxed implementation (looks "colour" and "mono")
% ---------------------------------------------------------------------------
%  Selected by style/theorems.tex, which documents the commands. Colours and
%  frames come from style/palette.tex; the paper look lives in
%  style/theorems-paper.tex and must define exactly the same commands.
% ===========================================================================

% ---------------------------------------------------------------------------
%  Box definitions
% ---------------------------------------------------------------------------
\newtcbtheorem[number within=section]{mydefinition}{Definition}{projectdefnbox}{defn}

\newtcbtheorem[use counter from=mydefinition]{mytheorem}{Theorem}{projectthmbox}{thm}

\newtcbtheorem[use counter from=mydefinition]{mylemma}{Lemma}{projectlembox}{lem}

\newtcbtheorem[use counter from=mydefinition]{mycorollary}{Corollary}{projectcorbox}{cor}

\newtcbtheorem[use counter from=mydefinition]{myproposition}{Proposition}{projectpropbox}{prop}

\newtcbtheorem[use counter from=mydefinition]{myclaim}{Claim}{projectclmbox}{clm}

\newtcbtheorem[use counter from=mydefinition]{myfact}{Fact}{projectfactbox}{fact}

% ---------------------------------------------------------------------------
%  Proof blocks (side rule matching the parent box, see style/palette.tex)
% ---------------------------------------------------------------------------
\NewDocumentEnvironment{projectrulebox}{m}{
    {\noindent{\it \textbf{Proof:}}}
    \tcolorbox[blanker,breakable,left=5mm,parbox=false,
        before upper={\parindent0pt},
        after skip=10pt,
        borderline west={1mm}{0pt}{#1}]
}{
    \endtcolorbox
}

\newenvironment{thmpf}{
    \begin{projectrulebox}{projectthmrule}
}{
    \textcolor{projectthmrule}{\hbox{}\nobreak\hfill$\blacksquare$}
    \end{projectrulebox}
}

\newenvironment{lempf}{
    \begin{projectrulebox}{projectlemrule}
}{
    \textcolor{projectlemrule}{\hbox{}\nobreak\hfill$\blacksquare$}
    \end{projectrulebox}
}

\newenvironment{corpf}{
    \begin{projectrulebox}{projectcorrule}
}{
    \textcolor{projectcorrule}{\hbox{}\nobreak\hfill$\blacksquare$}
    \end{projectrulebox}
}

\newenvironment{proppf}{
    \begin{projectrulebox}{projectproprule}
}{
    \textcolor{projectproprule}{\hbox{}\nobreak\hfill$\blacksquare$}
    \end{projectrulebox}
}

\newenvironment{clmpf}{
    \begin{projectrulebox}{projectclmrule}
}{
    \textcolor{projectclmrule}{\hbox{}\nobreak\hfill$\blacksquare$}
    \end{projectrulebox}
}

% Neutral standalone proof, also usable with an own heading:
%   \begin{projectproof}[Sketch of proof.] ... \end{projectproof}
\newenvironment{projectproof}[1][Proof.]{
    \par
    \noindent\textit{\textbf{#1}}\par
    \begin{tcolorbox}[
        blanker,
        breakable,
        left=5mm,
        parbox=false,
        before upper={\parindent0pt},
        after skip=10pt,
        borderline west={1mm}{0pt}{projectproofrule}
    ]
}{
    \hfill$\blacksquare$
    \end{tcolorbox}
}

\NewDocumentCommand{\pf}{+m}{
    \begin{projectproof}
        #1
    \end{projectproof}
}

% ---------------------------------------------------------------------------
%  Definition
% ---------------------------------------------------------------------------
\NewDocumentCommand{\defn}{o m+m}{
    \IfNoValueTF{#1}{\index{#2}}{%
        \def\defn@key{#1}\def\defn@nosig{-}%
        \ifx\defn@key\defn@nosig\else\index{#1}\fi%
    }%
    \begin{mydefinition}{#2}{}
        #3
    \end{mydefinition}
}

\NewDocumentCommand{\defnr}{o mm+m}{
    \IfNoValueTF{#1}{\index{#2}}{%
        \def\defn@key{#1}\def\defn@nosig{-}%
        \ifx\defn@key\defn@nosig\else\index{#1}\fi%
    }%
    \begin{mydefinition}{#2}{#3}
        #4
    \end{mydefinition}
}

% ---------------------------------------------------------------------------
%  Theorem
% ---------------------------------------------------------------------------
\NewDocumentCommand{\thm}{m+m}{
    \begin{mytheorem}{#1}{}
        #2
    \end{mytheorem}
}

\NewDocumentCommand{\thmr}{mm+m}{
    \begin{mytheorem}{#1}{#2}
        #3
    \end{mytheorem}
}

\NewDocumentCommand{\thmp}{m+m+m}{
    \begin{mytheorem}{#1}{}
        #2
    \end{mytheorem}

    \begin{thmpf}
        #3
    \end{thmpf}
}

\NewDocumentCommand{\thmpr}{mm+m+m}{
    \begin{mytheorem}{#1}{#2}
        #3
    \end{mytheorem}

    \begin{thmpf}
        #4
    \end{thmpf}
}

% ---------------------------------------------------------------------------
%  Lemma
% ---------------------------------------------------------------------------
\NewDocumentCommand{\lem}{m+m}{
    \begin{mylemma}{#1}{}
        #2
    \end{mylemma}
}

\NewDocumentCommand{\lemr}{mm+m}{
    \begin{mylemma}{#1}{#2}
        #3
    \end{mylemma}
}

\NewDocumentCommand{\lemp}{m+m+m}{
    \begin{mylemma}{#1}{}
        #2
    \end{mylemma}

    \begin{lempf}
        #3
    \end{lempf}
}

\NewDocumentCommand{\lempr}{mm+m+m}{
    \begin{mylemma}{#1}{#2}
        #3
    \end{mylemma}

    \begin{lempf}
        #4
    \end{lempf}
}

% ---------------------------------------------------------------------------
%  Corollary (no title)
% ---------------------------------------------------------------------------
\NewDocumentCommand{\cor}{+m}{
    \begin{mycorollary}{}{}
        #1
    \end{mycorollary}
}

\NewDocumentCommand{\corr}{m+m}{
    \begin{mycorollary}{}{#1}
        #2
    \end{mycorollary}
}

\NewDocumentCommand{\corp}{+m+m}{
    \begin{mycorollary}{}{}
        #1
    \end{mycorollary}

    \begin{corpf}
        #2
    \end{corpf}
}

% ---------------------------------------------------------------------------
%  Proposition (no title)
% ---------------------------------------------------------------------------
\NewDocumentCommand{\prop}{+m}{
    \begin{myproposition}{}{}
        #1
    \end{myproposition}
}

\NewDocumentCommand{\propr}{m+m}{
    \begin{myproposition}{}{#1}
        #2
    \end{myproposition}
}

\NewDocumentCommand{\propp}{+m+m}{
    \begin{myproposition}{}{}
        #1
    \end{myproposition}

    \begin{proppf}
        #2
    \end{proppf}
}

% ---------------------------------------------------------------------------
%  Claim (unnumbered)
% ---------------------------------------------------------------------------
\NewDocumentCommand{\clm}{m+m}{
    \begin{myclaim*}{#1}{}
        #2
    \end{myclaim*}
}

\NewDocumentCommand{\clmp}{m+m+m}{
    \begin{myclaim*}{#1}{}
        #2
    \end{myclaim*}

    \begin{clmpf}
        #3
    \end{clmpf}
}

% ---------------------------------------------------------------------------
%  Fact
% ---------------------------------------------------------------------------
\NewDocumentCommand{\fact}{+m}{
    \begin{myfact}{}{}
        #1
    \end{myfact}
}

% ---------------------------------------------------------------------------
%  Side-rule blocks: example, remark, unnumbered fact
% ---------------------------------------------------------------------------
\NewDocumentEnvironment{projectsideblock}{m}{
    \par
    \vspace{5pt}
    \noindent\textbf{#1}
    \begin{tcolorbox}[
        blanker,
        breakable,
        left=5mm,
        parbox=false,
        before upper={\parindent0pt},
        after skip=10pt,
        borderline west={1mm}{0pt}{projectsiderule}
    ]
}{
    \end{tcolorbox}
    \vspace{5pt}
}

\newenvironment{example}{
    \begin{projectsideblock}{Example.}
}{
    \end{projectsideblock}
}

\NewDocumentCommand{\ex}{+m}{
    \begin{example}
        #1
    \end{example}
}

\newenvironment{remark}{
    \begin{projectsideblock}{Remark.}
}{
    \end{projectsideblock}
}

\NewDocumentCommand{\rmkb}{+m}{
    \begin{remark}
        #1
    \end{remark}
}

% Inline remark: a short italic blue aside inside running text
\NewDocumentCommand{\rmk}{+m}{
    {\itshape\color{projectrmk}#1}
}

% Unnumbered fact, styled like a remark (the numbered version is \fact)
\newenvironment{factenv}{
    \begin{projectsideblock}{Fact.}
}{
    \end{projectsideblock}
}

\NewDocumentCommand{\factb}{+m}{
    \begin{factenv}
        #1
    \end{factenv}
}
